The following schedules improve the \emph{combined} frontier, not just any
single supplied report. We reuse the supplied asymmetric pair-load lemma,
thin precision gaps, exact Hermite taper, and all-level full-block Cauchy
factorization. The new ingredients here are a finite succession of grouping
alphabets, a carefully allocated geometric tail, and, for low-degree inputs,
allowing generated packet degrees to grow while charging their actual size.
No optimality claim is made for the schedules.

\subsection{The local bound and what is being compared}
Put $\lambda=\lceil\log_2s\rceil$. At depth $d$, let $u_d$ be the retained
valuation range, $a_d$ the packet exponent, $M_d$ the active edge count, and
$r_d$ the grouping alphabet. The compatible supplied ingredients give
\begin{align}
 M_{d+1}&\le (2/r_d)(1+o(1))M_d,\label{eq:ms-child}\\
 \log_2 N_d&\le (K(r_d)+o(1))M_d/\sqrt{u_d},\qquad
 K(r)=\frac{\pi}{\ln2}\sqrt{2(1-1/r)}.\label{eq:ms-local}
\end{align}
All error terms below are uniform over active calls. To see the coefficient,
write $D=2^{a_d-i}$. The precision envelope is $(6+o(1))M_dD/u_d$;
the full-block mode-family envelope is
$(2(1-1/r_d)+o(1))M_d/D$. The joint weights have slope four:
$w_j=a_d+1-i+4j-2s_2(j)$. Their partition product is bounded above by
that for $4j+1-2s_2(j)$, whose logarithm is
$(\pi^2+o(1))/(24\theta)$ per family. Optimizing the positive real
parameter $\theta$ yields~\eqref{eq:ms-local}. Dropping the positive
intercept is an upper bound; no assertion about a sharp joint limit of the
intercept and saddle is needed.

For comparison, the supplied asymmetric schedule has leading coefficient
$76.9089949516\ldots$ with slope-three modes. Simply substituting the
supplied joint modes gives $66.6051434076\ldots$, without any new schedule.
The asymmetric report also gives an excluded-boundary constant
$76.4327736858\ldots$; multiplying this stronger comparison by
$\sqrt{3/4}$ gives $66.1927236936\ldots$. More explicitly, its limiting
$r=16$ root budget is $1/80$ and the remaining budget is $1/80$; optimally
allocating the latter geometrically gives
\[
 \frac{K(16)}{\sqrt{1/80}}+
 \frac{(1/8)K(16)}{\sqrt{1/80}(1-(1/8)^{2/3})^{3/2}}
 =66.1927236936\ldots.
\]
This boundary is an infimum of admissible fixed-slack schedules, not an
attained finite parameter choice. The new bound below is smaller even
than it. A stronger retuned comparison allows $r_0=17$ at the root and $r=13$
thereafter. With $V=1/(20\log_2 17)$, $B=1/(10\log_2 13)-V$, its
excluded-boundary value is
\[
 \frac{K(17)}{\sqrt V}+
 \frac{(2/17)K(13)}{\sqrt B(1-(2/13)^{2/3})^{3/2}}
 =66.1204754600\ldots.
\]
The new theorem also improves on this more generous comparison. We do not
claim this two-stage value is globally minimal or that all possible
recombinations of the supplied ideas have been optimized.
These comparisons concern upper-bound certificates, not lower bounds on
previous algorithms.

\begin{theorem}[Generic polynomial-space coefficient below 64]
\label{thm:multistage-generic}
On the original normalized interface with at most $s$ edges, maximum degree
at most $\lceil s^{9/10}\rceil$, $P_0,\ell\le s$, polynomial description
length dependence, and original entry allowance
$\operatorname{poly}(s+\mathcal L)2^{A s^{9/10}}$ for fixed $A$, the unchanged
source identities together with the supplied local lemmas admit contraction
in polynomial space and time
\[
 (s+\mathcal L)^{O(1)}
 2^{(63.7573555464258\ldots+o(1))s/\sqrt\lambda}.
\]
In particular the coefficient can be replaced by $64$ for sufficiently
large $s$. This is a relative contraction theorem, not an independent
verification of the whole source or a coefficient directly on SAT variables.
\end{theorem}

\subsection{An entirely rational schedule}
Use $r_0,\ldots,r_5=(13,14,8,6,5,4)$ and $r_d=3$ for $d\ge6$.
Define the rational cumulative leading budgets
\[
 (A_0,\ldots,A_5)
 =(.01351,.02626,.03333,.03868,.04306,.04999),
 \qquad A_*=.06308,\qquad q=\frac{763}{1000}.
\]
Every displayed terminating decimal here denotes its exact rational value.
Put $A_{-1}=0$, $v_d=A_d-A_{d-1}$ for $0\le d\le5$, and
\[
 v_{6+j}=(A_*-A_5)(1-q)q^j\quad(j\ge0),\qquad
 u_d=\lfloor\lambda v_d\rfloor.
\]
Set
\[
 g=\lceil\log_2\lambda\rceil^2,\qquad
 a_d=\sum_{j=0}^d u_j+(d+1)g,\qquad H=\lceil s/\lambda\rceil.
\]
Below a finite threshold, or if the required integer inequalities fail,
use direct contraction. Thus no small input is forced into a degenerate
zero-range taper. The fixed rational slack is deliberate: the theorem does
not rely on implementing an excluded boundary or using floating-point logs.

\begin{proof}[Proof of Theorem~\ref{thm:multistage-generic}]
We first establish uniform applicability without assuming a recursion depth.
Consider every hypothetical depth
$0\le d\le D_*:=\lceil\log_{10/7}\lambda\rceil+10$ and every active
$M>H$. In this range $(d+1)g=O((\log\lambda)^3)=o(\lambda)$.
For the six finite stages, $a_d\le A_d\lambda+o(\lambda)$; in the tail,
$a_d\le A_*\lambda+o(\lambda)$. Every $u_d$ tends to infinity uniformly:
indeed $q>7/10$ gives
\[
 \min_{d\le D_*}u_d\ge\lambda^{\eta+o(1)},\qquad
 \eta=1+\frac{\ln q}{\ln(10/7)}>0.
\]
Floors are consequently a uniform relative $o(1)$ perturbation.

At the root, symmetric packing needs
$r_0^{2a_0}s^{9/10}\operatorname{poly}(\lambda)=o(M)$.
The leading inequality is
$20A_0\log_2(13)<1$. Below the root, the supplied child-side invariant
is $R_d\le\lceil s^{9/10}\rceil$ and $C_d\le2^{a_{d-1}+1}$.
The asymmetric packing lemma requires
\[
 Q_d R_d/M=o(1),\quad
 Q_d^2 C_d\lambda^2/M=o(1),\quad
 2^{a_d}Q_d^2/M=o(1),\qquad Q_d=r_d^{a_d}.
\]
For the row term it suffices that $10A_d\log_2r_d<1$ at finite stages
$d\ge1$, and $10A_*\log_2 3<1$ in the tail. Every one of these
inequalities is strict and checked by integer powers: if $A=n/m$,
check $13^{20n}<2^m$ at the root and $r^{10n}<2^m$ thereafter.
No approximate logarithm is used in these checks. The other tail exponents
are smaller than $A_*+1/5=.26308<1$. The root rounding exponent is less
than $A_0+1/10<1$. All logarithmic factors, $s/M\le\lambda$, ceilings,
and the additive precision gaps are absorbed by the fixed positive margins.
Also $2^{a_d+1}<s^{9/10}$ eventually, so no created tensor violates the
high-degree premise. Uniform near-unit padding follows.

For sufficiently large $s$, each hypothetical child shrinks by at least a
factor $7/10$, since the largest limiting ratio is $2/3$.
An active recursion path cannot reach $D_*$. This closes the induction and
validates the hypothetical-depth argument. More sharply, the padding bound
is $1+O(1/\lambda)+s^{-\epsilon+o(1)}$ for a fixed $\epsilon>0$;
its product over $O(\log\lambda)$ depths is $1+o(1)$. Therefore, with
$m_0=1$ and $m_d=\prod_{j<d}(2/r_j)$,
\[
 M_d\le(1+o(1))m_ds.
\]
The inherited precision is at most
$s\lambda^{O(1)}2^{a_{d-1}}$. Since every retained $i<u_d$ has
$2^{a_d-i}\ge2^{a_{d-1}+g+1}$, its ratio to
$M_d2^{a_d-i}/u_d$ is at most $\lambda^{O(1)}2^{-g}=o(1)$.
The evaluation loss is negligible by the same estimate. Thus the supplied
exact taper yields the uniform six-bit envelope used in~\eqref{eq:ms-local}.
All old tensor descriptions retain their creation precisions and canonical
representatives; the schedule changes none of these algebraic conventions.

Summing~\eqref{eq:ms-local} along a worst path gives the coefficient
\begin{equation}\label{eq:ms-constant}
 C=\sum_{d=0}^{5}\frac{m_dK(r_d)}{\sqrt{v_d}}
 +\frac{m_6K(3)}{\sqrt{(A_*-A_5)(1-q)}
                   (1-(2/3)/\sqrt q)}
 =63.7573555464258\ldots.
\end{equation}
The denominator is positive because $q>(2/3)^2$.
The six finite contributions are, respectively,
$52.98228468$, $8.41547470$, $1.56719246$, $0.43954213$,
$0.15865534$, and $0.04885070$; the infinite tail contributes
$0.14535554$. These decimals explain the allocation but are not proof inputs.
The accompanying exact-rational checker bounds square roots by integer
squares, $\pi$ by alternating Machin-series bounds, and $\ln2$ by its
positive $2\operatorname{atanh}(1/3)$ series with an explicit tail bound.
It certifies $C<64$.

There are polynomially many interpolation nodes, values, and child-group
choices per call. Across $O(\log\lambda)$ depths their logarithmic overhead
is $O(\lambda\log\lambda)=o(s/\sqrt\lambda)$. Direct leaves cost
$2^{O(s/\lambda)}$. Entry work is bounded by the original
$2^{O(s^{9/10})}$ allowance and packet work
$\exp(O(2^{a_d}\log s))=\exp(s^{A_*+o(1)})$; both are lower order.
All precision, arithmetic, label-set, and deformation-degree parameters
remain polynomial in $s$; in particular $M(2L-1)<s^4$ eventually.
Depth-first accumulation with immutable descriptions gives polynomial space.
The dependence on description length stays a fixed polynomial because
branching depends on $s$, not that length. This proves the theorem.
\end{proof}

\subsection{Tracking generated degrees improves the low-degree coefficient}
The supplied low-degree theorem holds all generated degrees below
$s^{1/12}$. This is sufficient but unnecessary: a later call can pay for the
actual degree of its immediately preceding packet scale instead. The fact
that all retained vertices move to one side does not create new incident
ports. Explicitly, the fresh side has $C_d\le 2L_{d-1}$, while the
retained side has
$R_d\le\max\{\lceil s^{1/12}\rceil,2L_{d-2},2L_{d-3},\ldots\}$;
in particular both sides are bounded by
$s^{\max\{1/12,A_{d-1}\}+o(1)}$.
Consequently, when packet exponents are nondecreasing and exceed
the initial degree exponent, \emph{both} child sides have maximum degree
at most $2^{a_{d-1}+1}$: retained older packets are no larger than the last
ones, and original vertices are smaller still. Symmetric matching packing
already suffices for the following theorem; asymmetric packing is not needed
in this smaller-degree branch.

\begin{theorem}[Low-degree polynomial-space coefficient below 24]
\label{thm:multistage-lowdegree}
Under the same normalized interface, but with root maximum degree at most
$\lceil s^{1/12}\rceil$, there is polynomial-space contraction in time
\[
 (s+\mathcal L)^{O(1)}
 2^{(23.6661688071455\ldots+o(1))s/\sqrt\lambda}.
\]
Thus the coefficient $24$ is eventually valid. This premise holds for the
direct bounded-degree sparse-CNF network, not for arbitrary dense incidence
networks merely because they arise from SAT.
\end{theorem}

\begin{proof}
Use $r_0,\ldots,r_5=(35,11,7,6,5,4)$, then $r_d=3$, with
\[
 (A_0,\ldots,A_5)
 =(.08881,.12627,.15116,.16206,.17717,.19999),
 \qquad A_*=.2398.
\]
Use exactly the preceding rational $q$, definitions of $v_d,u_d,a_d$, and
leaf threshold. Here $A_0>1/12$ and all $A_d$ increase strictly, so after
the root the maximum degree is at most $2^{a_{d-1}+1}$ as just explained.
Every finite stage and the tail satisfy
\[
 A(1+2\log_2r)<1.
\]
For $A=n/m$, the exact integer check is $r^{2n}<2^{m-n}$.
These inequalities control both rounding $LQ^2$ and the packing term
$Q^2\Delta$: below the root its leading exponent is
$2A_d\log_2r_d+A_{d-1}<A_d(1+2\log_2r_d)$, and at the root
$2A_0\log_2r_0+1/12<A_0(1+2\log_2r_0)$.
All inequalities have fixed strict margins. The same noncircular depth,
uniform floor, precision, and path estimates apply. The maximum packet
exponent is now $.2398\lambda+o(\lambda)$, still smaller than $.9\lambda$
and well within polynomial arithmetic and lower-order entry work.

Formula~\eqref{eq:ms-constant} with these parameters is
$23.6661688071455\ldots<24$, certified by the same exact arithmetic.
The six finite contributions are $21.19892854$, $1.80435144$,
$0.39079845$, $0.16636684$, $0.04614898$, and $0.01454394$; the
geometric tail adds $0.04503062$.
For comparison, combining joint modes with the supplied low-degree
$31.0641320630\ldots$ bound gives $26.9023275131\ldots$.
Again the new theorem improves the combined bound, not just its slope-three
component. The source SAT reduction constants are still needed to convert
any of these wire-level coefficients to a coefficient on $n$.
\end{proof}

\subsection{What these schedules do not establish}
The generic and low-degree certificates preserve the
$s/\sqrt{\log s}$ scale. General-density reductions still carry their
$O(n\log\log n)$ wire budget; these proofs neither remove that factor nor
provide a smaller unrestricted logarithmic order. Neither coefficient is a
measured running time or a formalized end-to-end SAT theorem. The auxiliary
search explored a restricted saturated-milestone schedule family; its
numerical candidates are not optimality certificates. The rational schedules
and all-size arguments above stand independently of that search.
